\section{\label{I_1_B}Son adaptation à l'archivistique et au contexte muséal}

Le modèle \gls{FRBR} est un modèle conçu par et pour les bibliothèques. À ce jour, les bibliothèques n'utilisent plus seulement le modèle FRBR mais la somme du FRBR et de ses modèles étendus : le \gls{FRAD} (Fonctionnalités requises des données d'autorité) pour la description des autorités, et le FRSAD (Fonctionnalités requises des données d'autorité matière) pour la description des sujets.\footcite{zotero-item-651}. La somme de ces modèles constitue le modèle \gls{IFLA-LRM}, successeur du \gls{FRBR}. Le modèle \gls{FRBR} est toutefois largement utilisé dans le domaine de la gestion des collections muséales. Or, sa qualité de modèle ne prétend pas normaliser la description des données, seulement d’en cadrer l’affichage. De fait, les entités proposées par le modèle \gls{FRBR} peuvent paraître génériques, et doivent être soumises à une redéfinition à chaque implémentation selon le contexte donné \footcite{boeuf2008}. Cette flexibilité ne doit toutefois pas nuire à l'objectif principal, celui de faciliter la navigation de l'utilisateur, en explicitant les “interrelations" entres les notices et les données produites \footcite{boeuf2008}. Nous étudierons dans les chapitres suivants une double adaptation, celle du modèle \gls{FRBR} à d’autres types de documents : objets de collections et archives, et son adaptation à différents contextes institutionnels pour une même nature de documents, celle des archives audiovisuelles.  

%Analysons premièrement son adaptation au milieu archivistique. La difficulté principale de l'adaptation du modèle FRBR aux collections d'archives est que le niveau Œuvre ne peut contenir le plus haut niveau archivistique : le fonds d'archives. En effet, un fonds d'archive ne peut être considéré systématiquement comme une œuvre, puisque dans la plupart des cas, il s'agit plutôt d'une agrégation naturelle d'éléments quotidiens produits par une personne morale ou physique. Ainsi, un fonds d'archives ne peut être intégré à la définition d'une œuvre comme création intellectuelle et artistique distincte \footcite{fick2015}. Un fonds d'archives correspondrait davantage à un agrégat d'œuvres ou de travaux, niveau intermédiaire décrit par le modèle FRBR : "un ensemble de documents privés classés par des archives en un seul fonds"(traduction). \footcite{fick2015}. Le niveau manifestation est obsolète car les archives ne sont pas publiées ?  

%Développer un peu plus sur l’adaptation du modèle aux archives : qu’en est-il des autres niveaux ? À quoi ressemble un catalogue d’une institution qui utilise le FRBR pour les archives ? (À relier avec le sujet) 

Le modèle de données \gls{FRBR} a été adapté au milieu muséal à travers une version orientée objet, le \gls{FRBRoo}, extension du \gls{CIDOC-CRM} (object-oriented Conceptual Reference Model), l'ontologie de référence dédiée au patrimoine culturel et muséal \footcite{gillis2015}. Il s’agit d’un modèle conceptuel général visant l’interopérabilité entre les données provenant d’institutions diverses dont les musées, centres d’archives et bibliothèques.  

Ce dernier modélise une approche descriptive centrée sur les évènements entourant l’objet décrit. Ainsi, les entités bibliographiques ne sont plus seulement décrites par des attributs, mais par des évènements spatio-temporels qui les créent et les font évoluer selon différents contextes \footcite{dialogue2017} Cette adaptation du \gls{FRBR} est particulièrement pertinente pour la description de collections de musées composites et de différentes natures, pouvant englober à la fois la description de tableaux, sculptures, ouvrages, specimens, films, etc. Tandis que le \gls{FRBR} est modélisé à partir de la nature d’un document bibliographique le FRBRoo replace la description dans un contexte plus large, celui d’évènements situés dans l’espace et dans le temps, évènements qui peuvent agir sur n’importe quel type d’objet.  

Le \gls{FRBRoo} est un modèle intéressant pour la gestion des archives audiovisuelles. En effet, les \gls{Archivesaudiovisuelles} possèdent un statut hybride au sein d’institutions patrimoniales, relevant à la fois de la création artistique, de l’œuvre et d’un archivage du réel, d’un témoignage \footcite{zotero-item-669}. Le \gls{FRBRoo} propose deux entités pertinentes, qui pourraient permettre de distinguer l’évolution de ce double statut : l’œuvre (Individual work), et la captation ou l’enregistrement (Recording work). Par exemple, l’étape du tournage serait documentée par l’entité Captation ainsi que l'ensemble des évènements et des personnes qui y sont liés, alors que le montage serait documenté par l’entité Œuvre ainsi que les évènements et les personnes qui y sont liés. Le FRBRoo conserve ensuite les entités prévues par le FRBR : l’Expression, la Manifestation, l’Item.  

Néanmoins, si le FRBRoo constitue une adaptation pertinente du FRBR pour le milieu muséal, sa structure trop large ne permet pas de rendre compte la nature propre du médium cinématographique, qui fonde la gestion des collections audiovisuelles. Ainsi, les institutions préservatrices de collections audiovisuelles se tournent davantage vers des adaptations métier spécifiques, comme celle de la fédération des archives du film (FIAF). La FIAF propose en effet une application du FRBR spécifique aux usages métier de description des archives audiovisuelles, au sein d’un manuel de catalogage des images animées. 
